\documentclass[a4paper,12pt]{article}
\usepackage[slovak]{babel}
\usepackage[utf8]{inputenc}
\usepackage[T1]{fontenc}
\usepackage{graphicx}
\usepackage{hyperref}
\usepackage{enumitem}
\usepackage{titlesec}
\usepackage{geometry}
\geometry{margin=2.5cm}

\titleformat{\section}{\large\bfseries}{\thesection.}{0.5em}{}
\titleformat{\subsection}{\normalsize\bfseries}{\thesubsection.}{0.5em}{}
\setlist[itemize]{noitemsep, topsep=0pt}



\begin{document}

%---------------------------------
% Titulná strana
%---------------------------------
\begin{titlepage}
    \centering
    \vspace*{3cm}
    {\Huge \textbf{Tvorba užívateľských rozhraní}}\\[1em]
    {\Large \textbf{Návrh aplikácie}}\\[5em]
    
    {\Large \textbf{Tím: Chleba}}\\[2em]
    {\large
        \textbf{Šimon Schön (xschons00)}\\
        \textit{Kamil Jakubčák (xjakubk00)}\\
        \textit{Adrián Pitka Kester (xpitkaa00)}
    }\\[6em]
    
    {\large \today}\\[20em]
    
    {\large Fakulta informačných technológií, VUT v Brne}
    \vfill
\end{titlepage}

%---------------------------------
% Obsah
%---------------------------------
\tableofcontents
\newpage

%---------------------------------
% Text dokumentu
%---------------------------------
\section{Návrh témy}
\begin{itemize}
  \item \textbf{Názov aplikácie:} Kapuut \\
\end{itemize}

Naša aplikácia má slúžiť ako zábavno edukatívna hra, v podobe kvízu s viacerými hernými režimami. Aplikácia je určená pre jedného hráča v režime pre precvičovanie otázok a odpovedí, alebo viacerých hráčov na jednom zariadení, v režime súťažného kvízu.

\section{Používateľský prieskum a špecifikácia}

\subsection{Prieskum [xschons00]}
Cieľom prieskumu bolo zistiť, ako používatelia pristupujú k vedomostným hrám a čo zvyšuje ich motiváciu k opakovanému hraniu. Zamerali sme sa na študentov, ktorí si chcú zábavnou formou precvičovať vedomosti z rôznych oblastí.
\\

\textbf{Profil používateľa:}  
Študent vysokej školy (vek 22 rokov), ktorý rád hrá mobilné hry vo voľnom čase.
\\

\textbf{Otázky užívateľského prieskumu:\\}
\begin{itemize}
  \item Aké typy vedomostných hier hrávaš najčastejšie a pri akej príležitosti?  
  \item Čo ťa motivuje pokračovať v hraní takýchto hier?  
  \item Uprednostňuješ kvíz, kde hráš sám, alebo s viacerými hráčmi?  
  \item Aký typ otázok (textové, obrázkové, časovo obmedzené) považuješ za zábavný?  
  \item Ako dlho by mala trvať jedna hra, aby udržala tvoju pozornosť?  
  \item Chcel by si mať možnosť pridať vlastné otázky?  
  \item Aké grafické prvky by ti spríjemnili hru (napr. animácie, zvukové efekty)?
  \item Uprednostňuješ hry, ktoré nevyžadujú pripojenie na internet?\\
\end{itemize}

\textbf{Zistenia:}  
Používateľ preferuje jednoduché a intuitívne prostredie, ktoré nevyžaduje zložité menu. Dôležitá je pre neho zábavná forma prezentácie otázok (napr. otáčanie kolesa, farebná spätná väzba po odpovedi) a možnosť rýchleho striedania hráčov. Oceňuje, ak hra nevyžaduje sieťové pripojenie a ponúka rýchle kolá v trvaní niekoľkých minút.

\subsection{Prieskum [xjakubk00]}
\begin{itemize}
    \item \textbf{Profil užívateľa}
    \begin{itemize}
        \item Vek: 14–15 rokov
        \item Kontext: pripravuje sa na prijímacie skúšky na strednú školu
        \item Cieľ: precvičiť otázky z matematiky, slovenčiny a logiky zábavnou formou
    \end{itemize}

    \item \textbf{Otázky pre rozhovor}
    \begin{itemize}
        \item \textbf{1. Všeobecné zistenia}
        \begin{itemize}
            \item Ako sa zvyčajne pripravuješ na prijímačky alebo testy?
            \item Používaš nejaké aplikácie alebo weby na precvičovanie? Ktoré a prečo?
            \item Čo ti na týchto aplikáciách chýba alebo ťa naopak baví?
        \end{itemize}

        \item \textbf{2. Vzťah k herným prvkom}
        \begin{itemize}
            \item Bavilo by ťa, keby učenie obsahovalo súťažný režim (napr. 1v1 proti spolužiakovi)?
            \item Aký typ hier ťa motivuje k ďalšiemu hraniu – zbieranie mincí, odomykanie úrovní, rebríčky?
            \item Čo by ťa viac motivovalo k precvičovaniu – súťaž, odmeny alebo sledovanie vlastného pokroku?
        \end{itemize}

        \item \textbf{3. Funkcie aplikácie Kapuut}
        \begin{itemize}
            \item Keby si si mohol vyberať medzi rôznymi témami (napr. matematika, gramatika, logika), ako by si chcel, aby to bolo usporiadané?
            \item Ako by si chcel, aby fungovala „ruleta“ s náhodnou otázkou – skôr pre zábavu alebo ako výzva?
            \item Mal by PvP režim obsahovať chat alebo iba herné výsledky?
            \item Ako by sa podľa teba mali používať mince – len na skóre alebo aj na odomykanie nových kategórií?
            \item Chceš vidieť správnu odpoveď po každej otázke alebo až na konci kola?
            \item Vadilo by ti, keby otázky obsahovali obrázky, grafy alebo krátke videá?
        \end{itemize}

        \item \textbf{4. Užívateľské rozhranie a prehľad}
        \begin{itemize}
            \item Ako by mal vyzerať hlavný panel – chceš vidieť svoj celkový progres, posledné hry alebo tipy na témy?
            \item Chcel by si, aby aplikácia pôsobila skôr ako hra (farebná, animovaná) alebo ako seriózny vzdelávací nástroj?
            \item Používaš častejšie mobil alebo počítač na učenie?
        \end{itemize}
    \end{itemize}

    \item \textbf{Úloha pre testovanie prototypu}
    \begin{itemize}
        \item Predstav si, že si v aplikácii Kapuut.
        \item Vyber si tému, zatoč ruletou a odpovedz na otázku.
        \item Po skončení kola si kúp novú kategóriu za mince alebo nejakú kozmetiku.
        \item Popíš, čo ti fungovalo dobre a čo by si zlepšil.
    \end{itemize}

    \item \textbf{Zhrnutie požiadaviek užívateľa}
    \begin{itemize}
        \item Používateľ chce mať možnosť učiť sa zábavnou formou a zároveň súťažiť s ostatnými.
        \item Užívatelia požadujú \textbf{okamžitú spätnú väzbu} – zobrazenie správnej odpovede po každej otázke.
        \item Je potrebné zobraziť \textbf{štatistiky pokroku} a \textbf{celkové skóre}.
        \item Odmeňovací systém s \textbf{mincami} musí motivovať k ďalšiemu učeniu (napr. odomykanie nových tém).
        \item Rozhranie má byť jednoduché, farebné a prispôsobené pre mobilné zariadenia.
    \end{itemize}
\end{itemize}


\subsection{Prieskum [xpitka00]}
TODO

\subsection{Zhrnutie používateľských požiadaviek [tímové]}
Konsolidovaný prehľad z prieskumov všetkých členov – popis kľúčových požiadaviek na funkčnosť a GUI.

\section{Prieskum existujúcich riešení}
\subsection{Prieskum [xschons00]}
\begin{itemize}
  \item \textbf{Kahoot:} Jedna z najznámejšich vedomostných hrier, používaných na školách. Je určené hlavne pre sútažné kvízy, kde sa hráči pripájajú na jedno zariadenie, ktoré slúži ako server. Používatelia môžu vytvárať a zdielať vlastné kvízy.
  \item \textbf{Sporcle:} Stránka zastrešuje mnoho rôznuch vedomostných minihier. Dizajn je trochu zastaralý ale ponúka viacero typov kvízov rozdelených do kategórií podĺa témy alebo typu hry.
\end{itemize}

\subsection{Prieskum [xpitka00]}
TODO


\subsection{Prieskum [xjakubk00]}

\begin{itemize}
    \item \textbf{Aplikácia 1: Quizizz}
    \begin{itemize}
        \item \textbf{Popis:}  
        Quizizz je vzdelávacia platforma zameraná na interaktívne kvízy pre študentov základných a stredných škôl. Ponúka režimy ako živé súťaže, domáce úlohy a adaptívne opakovanie.
        \item \textbf{Výhody / inšpirácie:}
        \begin{itemize}
            \item Okamžitá spätná väzba – po každej otázke sa zobrazí správna odpoveď a krátke vysvetlenie, čo zvyšuje efektivitu učenia.
            \item Gamifikácia – body, rebríček a motivačné vizuály udržiavajú pozornosť žiakov.
            \item Široké možnosti správy otázok – učitelia aj študenti môžu vytvárať vlastné sady s obrázkami a časovým limitom.
        \end{itemize}
        \item \textbf{Nevýhody / slabiny:}
        \begin{itemize}
            \item Žiadny PvP režim 1v1 – súťaž prebieha proti celej triede, nie proti konkrétnemu súperovi.
            \item Chýba prvok náhody – neexistuje „ruleta“ ani iný herný mechanizmus na náhodný výber otázok.
            \item Obmedzená personalizácia – neumožňuje prispôsobenie obtiažnosti alebo témy podľa hráča.
        \end{itemize}
    \end{itemize}

    \item \textbf{Aplikácia 2: Blooket}
    \begin{itemize}
        \item \textbf{Popis:}  
        Blooket kombinuje kvízy s hernými módmi (napr. vedomostné bitky, hazardné hry alebo tímové hry). Je obľúbený medzi žiakmi ZŠ a SŠ pre svoju zábavnosť.
        \item \textbf{Výhody / inšpirácie:}
        \begin{itemize}
            \item Viac herných módov – ponúka súťažné aj kooperatívne režimy, ktoré podporujú opakované hranie.
            \item Motivačný systém odmien – hráči zbierajú virtuálne predmety („blooks“) za správne odpovede.
            \item Jednoduché vytváranie a zdieľanie sád – komunita zdieľa tisíce hotových kvízov.
        \end{itemize}
        \item \textbf{Nevýhody / slabiny:}
        \begin{itemize}
            \item Nepodporuje flashcards ani individuálne precvičovanie mimo herný režim.
            \item Chýba priama stratégia hráča – neumožňuje vybrať otázku pre súpera ani skryť typ otázky.
            \item Príliš detský dizajn – menej vhodný pre starších študentov SŠ, ktorí preferujú serióznejší vizuálny štýl.
        \end{itemize}
    \end{itemize}

    \item \textbf{Závery a ponaučenia pre návrh aplikácie Kapuut}
    \begin{itemize}
        \item Kombinácia PvP súťaže s učením – na rozdiel od oboch konkurentov bude Kapuut ponúkať skutočný duel 1v1, kde hráči volia otázky pre súpera.
        \item Ruleta pre náhodnosť a napätie – funkcia „roulette“ pridá hernú náhodu a adrenalín, ktorý ostatným vzdelávacím aplikáciám chýba.
        \item Integrácia flashcards a štatistík – používateľ môže precvičovať témy samostatne, sledovať svoj pokrok a pripraviť sa na PvP zápasy.
        \item Systém mincí ako odmeny – hráči získavajú mince za správne odpovede, ktoré môžu využiť na odomykanie nových tém alebo vizuálnych prvkov, čo podporuje dlhodobú motiváciu.
    \end{itemize}
\end{itemize}



\begin{itemize}
  \item \textbf{Názov aplikácie 1:} [popis, účel]
  \item \textbf{Výhody / inšpirácie:} 
  \begin{itemize}
    \item [1)] konkrétna výhoda
  \end{itemize}
  \item \textbf{Nevýhody / nedostatky:}
  \begin{itemize}
    \item [1)] konkrétny nedostatok
  \end{itemize}
\end{itemize}

\subsection{Zhrnutie záverov tímu}
Tím definuje 3–5 kľúčových poznatkov a inšpirácií pre návrh vlastnej aplikácie.

\section{Rozdelenie práce v tíme}
Popis, kto rieši ktorú časť GUI alebo funkcie. Vysvetlenie, ako každá časť realizuje interaktívnu manipuláciu s dátami.

\section{Návrh GUI}
\subsection{Popis návrhu GUI jednotlivých členov}
Každý člen vloží návrh svojej časti GUI s popismi.
\begin{itemize}
  \item \textbf{Obrázok / wireframe:} \includegraphics[width=\textwidth]{gui_navrh1.png}
  \item \textbf{Popis:} ako GUI rieši používateľské potreby
  \item \textbf{Interaktívne prvky:} prehľad manipulácií s dátami
\end{itemize}

\section{Výber technológií}
Pre vývoj aplikácie sme zvolili \textbf{Godot Engine} – open-source herný engine vhodný na 2D hry s možnosťou multiplatformového exportu. Hlavné dôvody tohto výberu sú:

\begin{itemize}
  \item \textbf{Jednoduchosť použitia:} Godot poskytuje vizuálne prostredie, v ktorom je možné vytvárať scény a prepojiť ich so skriptami bez potreby rozsiahlej konfigurácie.
  \item \textbf{Rýchly vývoj prototypov:} Engine umožňuje okamžité testovanie a ladenie hry priamo počas vývoja.
  \item \textbf{Multiplatformovosť:} Export pre Windows, Linux, Android a web (HTML5) bez dodatočných úprav.
  \item \textbf{Open-source a komunita:} Veľká podpora zo strany komunity a množstvo dostupných návodov.
\end{itemize}

\section{Návrh SW architektúry a API (spoločné)}
\subsection{Architektúra systému}
Popis logických častí aplikácie, prepojenie FE a BE, prípadne schéma.

\subsection{Návrh API}
\begin{itemize}
  \item \textbf{Príklad funkcie:} \texttt{GET /api/users}
  \item \textbf{Vstupy:} [popis]
  \item \textbf{Výstupy:} [popis]
  \item \textbf{Prepojenie s GUI:} ktorá časť GUI túto funkciu volá
\end{itemize}

\subsection{Dátové štruktúry}
Stručné popisy modelov (napr. User, Item, Group...) a ich vlastností.

\section{Záver}
Zhrnutie zistení, prehľad splnených bodov, prípadné otvorené otázky.

\end{document}
